\section{ALCANCES Y LÍMITES}

Los alcances enumeran lo que el trabajo entrega; los límites, lo que queda fuera de él.

\subsection{Alcances}

\begin{itemize}
    \item Análisis de los procesos organizativos del Desafío Bebras y de los sistemas utilizados por otros países miembros, tomados como referencia funcional.
    \item Diseño de la arquitectura, del modelo de datos y de los componentes del sistema.
    \item Gestión de contenido: edición por secciones del contenido de las páginas del sitio informativo, administración del menú de navegación y del pie de página, creación de páginas nuevas, gestión de imágenes y noticias, y publicación de los cambios en el sitio.
    \item Gestión de desafíos: banco de tareas reutilizable con distintos tipos de respuesta y su probador; creación, configuración y publicación de desafíos por categoría, y modo de práctica con tareas de ejercitación.
    \item Registro y administración de participantes: registro y aprobación de maestros y unidades educativas; creación de grupos con código de acceso, carga de la nómina de estudiantes e ingreso al desafío por código.
    \item Evaluación: rendición en línea con autoguardado, corrección automática de las respuestas, generación y publicación de resultados, y generación de formatos imprimibles.
    \item Validación mediante pruebas funcionales por módulo y mediante la revisión del supervisor, incorporando las observaciones que formule hasta obtener su conformidad.
    \item Entrega del sistema desplegado y probado al equipo organizador, conforme al cronograma establecido.
    \item Elaboración del documento final, que registra el análisis, el diseño, el desarrollo y las pruebas de cada módulo.
\end{itemize}

\subsection{Límites}

\begin{itemize}
    \item El sistema se orienta exclusivamente a la gestión del Desafío Bebras en Bolivia para estudiantes de nivel primario y secundario.
    \item La gestión de contenido se limita a la edición del contenido, la navegación y la estructura del sitio web informativo; no comprende su rediseño gráfico ni la creación de nuevas plantillas visuales.
    \item No se contempla la integración con sistemas del Ministerio de Educación ni con plataformas educativas nacionales existentes.
    \item El trabajo concluye con la entrega del sistema desplegado y validado. El acompañamiento técnico se limita al periodo de entrega previsto en el cronograma; no se contempla la operación ni el mantenimiento permanente del sistema en ediciones posteriores.
    \item La ejecución del Desafío Bebras Bolivia y su operación en aulas no forman parte del trabajo: la organización y realización del Desafío corresponden al equipo organizador.
    \item El cumplimiento del trabajo no está condicionado a que el sistema sea utilizado en la edición 2026 del Desafío.
\end{itemize}
